\section{Related Work}
\label{sec:related}

\paragraph{Separating reasoning from authority.} Keeping authority away from the model that reads untrusted data is established. The dual-LLM
pattern~\cite{willison}, design patterns for resistant agents~\cite{patterns} and CaMeL~\cite{camel}, in which an interpreter checks capability
policies on every tool call, constrain what injected text can cause. CaMeL's policies run in the agent's host, so code execution there defeats them;
ours run on a chain and in a signer outside it. The two compose. PACE~\cite{pace} is the closest system: a deterministic verifier signs a decision record
for each transaction and a smart account checks it, binding policy approval to on-chain execution, as our vault's operator role does for x402 payments.
PACE checks one transaction at a time and reports that a multi-step drain passes its guards; our vault carries a cumulative budget and call count, and our
accounting model is about exactly those sequences. PACE does not hide its limits or study what repeated requests reveal; we hide them, bound what the
decisions reveal and measure how little that bound leaves. Its evaluation of model behaviour is broader than ours. Mao et al.~\cite{mao2026} recommend
separating cognition from custody with scoped spending policies at the signing layer, of which this work is an instance, and the Cloud Security Alliance
note~\cite{csa2026} reaches the same advice. Acharya~\cite{acharya2025} proposes an agent wallet that checks a user-signed mandate, without an implementation
and without considering probing. We claim novelty only for the confidential, cumulative construction on x402's escrow and for the failure analysis: we did
not find the interaction between refunds and probing in any of these works. AgentDojo~\cite{agentdojo} measures whether attacks succeed, including in a
banking suite; we bound what a successful one costs.

\paragraph{Agent payments and delegated spending.} x402~\cite{x402} and Google's Agent Payments Protocol~\cite{ap2} define how an agent pays and how a
mandate is expressed. Smart-account permissions~\cite{erc4337,erc7715,erc7710} let a session key or delegation spend only within limits the account enforces,
which holds under A2 and, unlike our design, also holds when the signing key is compromised, since the account is the payer. Our plaintext comparator is that
ordinary case in policy terms; an account whose validation consults our vault's record would remove the payer key from the trusted base, and we have not built
it. x402's auth-capture scheme~\cite{x402authcapture} and Offer and Receipt extension~\cite{x402offer} give separate capture and signed terms, not delivery.

\paragraph{Confidential contracts and rollback.} Hawk~\cite{hawk} and Ekiden~\cite{ekiden} keep contract state private with cryptography and trusted hardware;
Zama's fhEVM~\cite{fhevm} and Inco Lightning~\cite{inco} expose encrypted types to ordinary contracts. Our enclave-held comparator anchors an enclave's state on
chain to stop rollback, as ROTE~\cite{rote} and Nimble~\cite{nimble} do with their own counters and logs.

\paragraph{Fair exchange.} Fair exchange needs a trusted third party~\cite{pagnia1999}; optimistic protocols involve it only on dispute~\cite{asokan1998}.
Hash-time-locked contracts~\cite{lightning} make a payment conditional on a preimage. Zero-knowledge contingent payment~\cite{maxwell2016,campanelli2017} and
FairSwap~\cite{fairswap} extend the condition to a property of the delivered data. Our protocol is a hash lock with a bond and a FairSwap-style dispute,
restricted to one predicate whose execution and calldata fit under 200{,}000 gas; the complete dispute transition is not built.

\paragraph{Information flow.} The separation of an untrusted requester from a small mediator follows the reference monitor concept~\cite{anderson}. The
probing attack is a failure of robust declassification~\cite{zdancewic2001}, and the attempt budget a quantitative bound in the sense of
Smith~\cite{smith2009}.

\section{Limitations}
\label{sec:limits}

The enclave ran on one rented testnet machine for one day: a campaign, one benchmark run and one wipe test. The matched baselines and the earlier
benchmarks used a local signer with a seed file, and Base cannot verify the enclave's attestation; it pins the enrollment address the enclave reports. The arguments of Section~\ref{sec:security} are by hand, and the contracts have not been audited. Inco's properties I1 to
I5 are assumed, and liveness depends on its covalidators. A compromised signer can sign a plain transfer, because the payer is a key; only a smart-account payer
would remove that. Requests on a goal are strictly sequential. The decision-channel theorem assumes no reveals on dropped branches, which we have not measured,
and its relation to the deployed transcript rests on the unproven assumptions of Section~\ref{sec:leakthm}. The injection study has one free model through one
gateway, five samples per arm, and an agent that is told the limits it is held to. The delivery protocol is specified and only its predicate is built; held services, chunked commitments for goods larger
than 4 KiB and any check of quality beyond a committed hash are future work. The system has not run on mainnet, and a vendor must accept the vault as the
escrow's operator.

\section{Conclusion}

An agent that reads what sellers write cannot be made trustworthy by inspecting that text, and it need not be. With payments collectable only through the
policy contract and a payer enrolled for the user's wallet, a fully compromised agent keeps a bounded spending authority over the life of a goal, and the
bound survives a lying RPC and dropped blocks. Keeping the limits secret is a smaller gain than it first appears: refunds made the private policy free to
question, an attempt budget restores a bound, and on realistic caps that bound leaves a probing vendor close to full knowledge after a handful of answers.
For agents that pay, the lasting design question is not how to hide a limit but how few questions to answer about it.

\paragraph{Artifact.} Source code, contracts, tests, run logs, raw benchmark records, the leakage computation and the injection study's harness and results
are at \url{https://github.com/adarshkr7/LedgerLighthouse}, tag \code{arxiv-v3}; its \code{ARTIFACT.md} maps each claim to its record. The run records are
in \code{docs/runs}, the paper's source in \code{paper/arxiv-v3}, the evaluation code and data in \code{paper/eval}, and the enclave's image digest, measurements and
policy in \code{services/signer} and \code{docs/runs}. Every data table in this paper is generated from those records by one script.
What reproduces offline: the Solidity suites, including the accounting invariants, which install a recorded snapshot of the escrow, collector and USDC at
Base Sepolia block 56{,}480{,}000; the TypeScript suites; the leakage computation; and the full flow on a local fork with a stand-in for Inco. What needs the
network: the fork suites that test against live state (they skip, and say so, without \code{ESCROW\_FORK\_RPC}), the live runs (Base Sepolia, test USDC, and
Inco's covalidators), the injection study (a model gateway), and anything about the enclave (an Oasis ROFL machine). Failed and stalled runs are kept beside
the successful ones.

\paragraph{Acknowledgements.} I thank Krishan Pratap Sharma (IIT Bombay) for reading, discussion and help with the formal framing.
